% !TeX program = xelatex
\documentclass[a4paper,12pt]{article}

% ===== 1. KHỔ GIẤY, LỀ TRANG, FONT =====
\usepackage[left=3cm,right=2cm,top=2cm,bottom=2cm,footskip=1cm]{geometry}
\usepackage{fontspec}

\IfFontExistsTF{Times New Roman}{
  \setmainfont{Times New Roman}
}{
  \IfFileExists{fonts/times.ttf}{
    \setmainfont{times.ttf}[
      Path=fonts/,
      BoldFont=timesbd.ttf,
      ItalicFont=timesi.ttf,
      BoldItalicFont=timesbi.ttf]
  }{
    \IfFontExistsTF{TeX Gyre Termes}{
      \setmainfont{TeX Gyre Termes}
      \PackageWarningNoLine{MauBaoCao}{Times New Roman unavailable; using TeX Gyre Termes}
    }{
      \setmainfont{DejaVu Serif}
      \PackageWarningNoLine{MauBaoCao}{Times New Roman unavailable; using DejaVu Serif}
    }
  }
}
\usepackage{indentfirst}
\usepackage{microtype}
\usepackage{graphicx}
\usepackage{float}
\usepackage{tikz}
\usepackage{array,tabularx,booktabs,longtable}
\usepackage{enumitem}
\usepackage{titlesec}
\usepackage{tocloft}
\usepackage{caption}
\usepackage{xurl}
\usepackage[unicode,hidelinks]{hyperref}
\urlstyle{same}

\renewcommand{\normalsize}{\fontsize{13}{19.5}\selectfont}
\AtBeginDocument{\normalsize}
\setlength{\parindent}{0.5cm}
\setlength{\parskip}{4pt}
\setlength{\emergencystretch}{2em}
\raggedbottom
\setlist{topsep=4pt,itemsep=2pt,parsep=0pt,leftmargin=*}

% ===== 2. BA CẤP TIÊU ĐỀ: 1 / 1.1 / 1.1.1 =====
\setcounter{secnumdepth}{3}
\setcounter{tocdepth}{3}
\renewcommand{\thesection}{\arabic{section}}
\renewcommand{\thesubsection}{\thesection.\arabic{subsection}}
\renewcommand{\thesubsubsection}{\thesubsection.\arabic{subsubsection}}
\titleformat{\section}[hang]
  {\normalfont\fontsize{16}{20}\selectfont\bfseries}
  {\thesection.}{0.6em}{}
\titleformat{\subsection}[hang]
  {\normalfont\fontsize{14}{18}\selectfont\bfseries}
  {\thesubsection}{0.6em}{}
\titleformat{\subsubsection}[hang]
  {\normalfont\fontsize{13}{17}\selectfont\bfseries}
  {\thesubsubsection}{0.6em}{}
\titlespacing*{\section}{0pt}{16pt}{8pt}
\titlespacing*{\subsection}{0pt}{12pt}{6pt}
\titlespacing*{\subsubsection}{0pt}{8pt}{4pt}

% ===== 3. MỤC LỤC VÀ SỐ TRANG =====
\renewcommand{\contentsname}{MỤC LỤC}
\renewcommand{\refname}{Tài liệu tham khảo}
\renewcommand{\cfttoctitlefont}{\fontsize{16}{20}\selectfont\bfseries}
\setlength{\cftbeforetoctitleskip}{0pt}
\setlength{\cftaftertoctitleskip}{18pt}
\renewcommand{\cftsecfont}{\normalfont\bfseries}
\renewcommand{\cftsecpagefont}{\normalfont\bfseries}
\renewcommand{\cftsubsecfont}{\normalfont}
\renewcommand{\cftsubsubsecfont}{\normalfont}
\renewcommand{\cftsecleader}{\cftdotfill{\cftdotsep}}
\cftsetindents{section}{0em}{2em}
\cftsetindents{subsection}{2em}{3em}
\cftsetindents{subsubsection}{5em}{4em}
\setlength{\cftbeforesecskip}{6pt}
\makeatletter
\def\ps@plain{%
  \let\@oddhead\@empty\let\@evenhead\@empty
  \def\@oddfoot{\hfil{\normalfont\fontsize{11}{13}\selectfont\thepage}\hfil}%
  \let\@evenfoot\@oddfoot}
\makeatother
\pagestyle{plain}

% ===== 4. BẢNG, HÌNH, CHÚ THÍCH =====
\renewcommand{\tablename}{Bảng}
\renewcommand{\figurename}{Hình}
\DeclareCaptionFont{reportcaption}{\fontsize{12}{15}\selectfont}
\captionsetup{font=reportcaption,labelfont=bf,labelsep=period,
  justification=centering,skip=6pt}
\newcolumntype{Y}{>{\raggedright\arraybackslash}X}
\newcommand{\bangfont}{\fontsize{12}{15}\selectfont}

% ===== 5. THÔNG TIN TRANG BÌA =====
\newcommand{\DonViChuQuan}{ĐẠI HỌC QUỐC GIA HÀ NỘI}
\newcommand{\TenTruong}{TRƯỜNG ĐẠI HỌC CÔNG NGHỆ}
\newcommand{\LoaiBaoCao}{BÁO CÁO TIẾN ĐỘ DỰ ÁN 1}
\newcommand{\TenMonHoc}{ĐIỆN TOÁN ĐÁM MÂY}
\newcommand{\TenDeTai}{SatiBot -- Quản lý và tương tác với các tài nguyên AWS bằng ngôn ngữ tự nhiên}
\newcommand{\GiangVien}{TS. Hoàng Xuân Tùng}
\newcommand{\SoNhom}{Nhóm 10}
\newcommand{\LopMonHoc}{INT3319 1}
\newcommand{\DiaDiemThoiGian}{Hà Nội, Tháng 10/2026}
\hypersetup{pdftitle={BÁO CÁO TIẾN ĐỘ DỰ ÁN 1 - SATIBOT},pdfauthor={Nhóm 10}}

\begin{document}

% ===== 6. TRANG BÌA =====
\begin{titlepage}
\thispagestyle{empty}
\begin{tikzpicture}[remember picture,overlay]
  \draw[line width=2pt]
    ([xshift=3cm,yshift=-2cm]current page.north west)
    rectangle ([xshift=-2cm,yshift=2cm]current page.south east);
  \draw[line width=0.5pt]
    ([xshift=3.12cm,yshift=-2.12cm]current page.north west)
    rectangle ([xshift=-2.12cm,yshift=2.12cm]current page.south east);
\end{tikzpicture}\par
\centering
\setlength{\parindent}{0pt}
\setlength{\parskip}{0pt}
\vspace*{0.8cm}
{\fontsize{14}{18}\selectfont\bfseries\DonViChuQuan\par}
\vspace{0.15cm}
{\fontsize{15}{20}\selectfont\bfseries\TenTruong\par}
\vspace{0.25cm}
\rule{5cm}{0.4pt}\par
\vspace{0.5cm}

\IfFileExists{Image/logo.png}{
  \includegraphics[width=3.8cm,height=3.8cm,keepaspectratio]{Image/logo.png}
}{
  \fbox{\parbox[c][3.4cm][c]{3.4cm}{\centering
    \fontsize{11}{14}\selectfont LOGO TRƯỜNG}}
}\par
\vspace{0.7cm}
{\fontsize{22}{27}\selectfont\bfseries\LoaiBaoCao\par}
\vspace{0.35cm}
\begin{minipage}{0.9\textwidth}\centering
  {\fontsize{15}{20}\selectfont\bfseries MÔN HỌC: \TenMonHoc\par}
\end{minipage}\par
\vspace{0.6cm}
{\fontsize{13}{17}\selectfont\bfseries ĐỀ TÀI\par}
\vspace{0.12cm}
\begin{minipage}{0.9\textwidth}\centering
  {\fontsize{16}{21}\selectfont\bfseries\TenDeTai\par}
\end{minipage}\par
\vspace{0.65cm}
{\fontsize{13}{19.5}\selectfont
\renewcommand{\arraystretch}{1.1}
\begin{tabular}{@{}>{\raggedright\arraybackslash\bfseries}p{5.2cm}>{\raggedright\arraybackslash}p{8cm}@{}}
  Giảng viên hướng dẫn: & \GiangVien\\
  Nhóm thực hiện: & \SoNhom\\
  Sinh viên thực hiện: & Cao Anh Minh \textbf{(Trưởng nhóm)}\\
                      & Trịnh Thiên Quang\\
                      & Phạm Đức Hùng\\
                      & Nguyễn Đức Huy\\
  Lớp môn học: &  \LopMonHoc
\end{tabular}\par}
\vfill
{\fontsize{13}{17}\selectfont\bfseries\DiaDiemThoiGian\par}
\vspace*{0.5cm}
\end{titlepage}

% ===== 7. MỤC LỤC TỰ ĐỘNG =====
\pagenumbering{roman}
\tableofcontents
\clearpage
\pagenumbering{arabic}

% ===== 8. NỘI DUNG BÁO CÁO =====

\section{TỔNG QUAN VỀ DỰ ÁN}

\subsection{Giới thiệu chung về tính năng dự án}

Hiện nay, Amazon Web Services (AWS) là một trong những nền tảng điện toán đám mây phổ biến nhất phục vụ học tập, nghiên cứu và triển khai ứng dụng. Tuy nhiên, việc khởi tạo và quản trị tài nguyên thông qua giao diện web (AWS Management Console) đòi hỏi nhiều thao tác thủ công, điều hướng qua nhiều màn hình cấu hình phức tạp và rất dễ nhầm lẫn. Trong khi đó, việc sử dụng giao diện dòng lệnh (AWS CLI) hoặc viết tập lệnh tự động lại yêu cầu ghi nhớ chính xác cú pháp câu lệnh, cấu trúc tệp cấu hình JSON và các tham số kỹ thuật. Đối với sinh viên, lập trình viên mới tiếp cận điện toán đám mây hoặc người quản lý kỹ thuật, các rào cản này làm chậm quá trình vận hành, đồng thời tiềm ẩn nguy cơ cấu hình sai gây mất an toàn thông tin hoặc phát sinh chi phí ngoài ý muốn.

Nhằm giải quyết vấn đề trên, nhóm thực hiện đề tài xây dựng SatiBot -- ứng dụng ChatOps hỗ trợ quản lý, cấu hình và giám sát tài nguyên AWS thông qua câu lệnh bằng ngôn ngữ tự nhiên tiếng Việt. Bằng cách kết hợp mô hình ngôn ngữ lớn (LLM) hỗ trợ cơ chế \textit{Function Calling} với thư viện AWS SDK (Boto3), SatiBot hoạt động như một trợ lý hỗ trợ thao tác hạ tầng. Người dùng chỉ cần gửi yêu cầu bằng tiếng Việt qua giao diện chat; hệ thống sẽ phân tích ý định, kiểm tra tính hợp lệ và độ an toàn của câu lệnh, sau đó gọi các API tương ứng của AWS và trả về kết quả thực thi kèm theo thông tin chi tiết, cảnh báo rủi ro hoặc hướng dẫn xử lý tiếp theo.

\subsection{Giới thiệu người dùng của dự án}

Hệ thống SatiBot hướng tới 03 nhóm người dùng chính với nhu cầu sử dụng thực tế như sau:

\begin{itemize}
  \item \textbf{Lập trình viên và kỹ sư vận hành (Developers / DevOps Engineers):} Cần khởi tạo nhanh môi trường thử nghiệm, cấu hình máy chủ, cập nhật mã nguồn hàm serverless hoặc xem log hệ thống trong quá trình phát triển ứng dụng. SatiBot giúp nhóm này thực hiện các tác vụ lặp đi lặp lại nhanh hơn ngay trên khung chat mà không cần chuyển đổi qua lại giữa nhiều trang trên AWS Console.
  \item \textbf{Sinh viên ngành Công nghệ thông tin / Điện toán đám mây:} Cần môi trường thực hành trực quan cho các môn học về mạng và hạ tầng đám mây. Hệ thống cung cấp cơ chế kiểm tra an toàn (guardrails), giúp sinh viên tránh các lỗi phổ biến như mở cổng mạng công khai nguy hiểm, quên tắt máy ảo hoặc chọn nhầm các gói cấu hình vượt quá phạm vi gói miễn phí (AWS Free Tier).
  \item \textbf{Quản lý kỹ thuật và chủ dự án (Tech Leads / Project Managers):} Cần theo dõi nhanh tình trạng hoạt động của tài nguyên, kiểm tra phân quyền (IAM Audit) và kiểm soát chi phí sử dụng hàng tháng (FinOps) bằng ngôn ngữ tự nhiên mà không cần thành thạo cú pháp dòng lệnh AWS CLI.
\end{itemize}

\subsection{Cách thức và quy trình người dùng sử dụng sản phẩm để đạt mục đích}

\begingroup
\setlist{topsep=2pt,itemsep=1pt,parsep=0pt,leftmargin=*}
\setlength{\parskip}{2pt}

Để đáp ứng nhu cầu quản trị hạ tầng cơ bản, SatiBot hỗ trợ 04 dịch vụ chính của AWS gồm: Amazon EC2, Amazon S3, AWS IAM và AWS Lambda. Dưới đây là chi tiết về mục đích và quy trình tương tác thực tế của người dùng đối với từng nhóm tính năng của 04 dịch vụ AWS nòng cốt.

\subsubsection{Nhóm kịch bản và quy trình quản trị Amazon EC2}

Dịch vụ máy chủ ảo (Amazon Elastic Compute Cloud - EC2) là hạ tầng phục vụ việc triển khai và chạy ứng dụng. SatiBot hỗ trợ 05 quy trình quản lý trạng thái và cấu hình máy ảo như sau:

\begin{enumerate}
  \item \textbf{Khởi tạo máy chủ chuẩn hóa và xử lý trùng lặp tên:}
  \begin{itemize}
    \item \textit{Mục đích:} Tạo nhanh máy ảo để thử nghiệm hoặc triển khai website mà không phải cấu hình thủ công nhiều bước, đồng thời xử lý trường hợp tên máy bị trùng.
    \item \textit{Quy trình người dùng:} 
    \begin{itemize}
      \item Nhắn yêu cầu tạo máy kèm mục đích sử dụng (ví dụ: "Tạo cho tôi một máy chủ Ubuntu để chạy web demo, mở cổng SSH và HTTP").
      \item Nếu tên máy đã tồn tại, hệ thống đưa ra 3 lựa chọn: tiếp tục sử dụng máy cũ, đổi tên máy mới hoặc xóa máy cũ để tạo lại.
      \item Nhận kết quả thành công gồm: Địa chỉ IP máy chủ, link tải file khóa SSH (.pem) và câu lệnh kết nối mẫu.
    \end{itemize}
  \end{itemize}

  \item \textbf{Kiểm soát chi phí và xử lý vượt hạn mức tài nguyên (Cost Guardrails \& Quota):}
  \begin{itemize}
    \item \textit{Mục đích:} Tránh việc chọn nhầm máy cấu hình cao gây phát sinh chi phí lớn, đồng thời hướng dẫn cách xử lý khi tài khoản chạm ngưỡng hạn mức tạo máy ảo (vCPU quota).
    \item \textit{Quy trình người dùng:}
    \begin{itemize}
      \item Nhắn yêu cầu loại máy mong muốn (ví dụ: \textit{''Tạo máy ảo c5.4xlarge 16 vCPU''} hoặc \textit{''Tạo thêm 1 máy ảo Ubuntu''}).
      \item Xem phản hồi từ hệ thống:
      \begin{itemize}
        \item \textit{Nếu yêu cầu máy cấu hình cao, chi phí lớn:} Hệ thống cảnh báo chi phí ước tính hàng tháng, tạm dừng thao tác và yêu cầu người dùng xác nhận nếu vẫn muốn tiếp tục khởi tạo.
        \item \textit{Nếu tài khoản chạm giới hạn vCPU:} Hệ thống giải thích nguyên nhân và liệt kê danh sách các máy ảo đang chạy để người dùng cân nhắc tắt bớt máy không dùng.
      \end{itemize}
      \item Người dùng chọn phương án: chuyển về gói miễn phí (Free Tier), tắt bớt máy cũ hoặc gửi yêu cầu nâng hạn mức lên AWS.
    \end{itemize}
  \end{itemize}

  \item \textbf{Tra cứu và điều khiển vòng đời máy chủ (Bật / Tắt / Khởi động lại / Xóa):}
  \begin{itemize}
    \item \textit{Mục đích:} Theo dõi danh sách máy ảo và điều khiển trạng thái bật/tắt kịp thời nhằm tránh lãng phí tài nguyên và kiểm soát chi phí vận hành.
    \item \textit{Quy trình người dùng:}
    \begin{itemize}
      \item Nhắn câu lệnh tra cứu hoặc điều khiển (ví dụ: \textit{''Liệt kê các máy ảo đang chạy''}, \textit{''Tắt máy ảo web-demo''}, \textit{''Khởi động lại máy i-0abc123''}, \textit{''Xóa máy ảo staging''}).
      \item Xem bảng thông tin gồm ID, tên máy, trạng thái, IP công khai và thời gian chạy; nhận thông báo trạng thái cập nhật (kèm IP mới khi bật lại máy).
      \item Với thao tác xóa: Hệ thống cảnh báo mất dữ liệu và yêu cầu người dùng gõ xác nhận trước khi thực hiện xóa máy.
    \end{itemize}
  \end{itemize}

  \item \textbf{Giám sát hiệu năng và phát hiện máy chủ lãng phí (FinOps \& CloudWatch):}
  \begin{itemize}
    \item \textit{Mục đích:} Theo dõi mức độ sử dụng tài nguyên để tắt các máy nhàn rỗi nhằm tiết kiệm chi phí, hoặc nhận biết máy bị quá tải để kịp thời nâng cấp.
    \item \textit{Quy trình người dùng:}
    \begin{itemize}
      \item Nhắn câu hỏi kiểm tra (ví dụ: \textit{''Kiểm tra CPU máy i-0abc123 trong 24 giờ qua''}).
      \item Nhận số liệu đánh giá chi tiết:
      \begin{itemize}
        \item \textit{Cảnh báo lãng phí:} Nếu máy chạy không tải kéo dài, hệ thống đưa ra gợi ý tắt máy để tiết kiệm chi phí.
        \item \textit{Cảnh báo nghẽn tải:} Nếu mức sử dụng CPU liên tục ở mức cao, hệ thống cảnh báo quá tải và gợi ý nâng cấp loại máy.
      \end{itemize}
      \item Người dùng nhắn đồng ý nếu muốn hệ thống tự động xử lý theo khuyến nghị.
    \end{itemize}
  \end{itemize}

  \item \textbf{Thiết lập tường lửa an toàn và hỗ trợ sự cố kết nối (SecOps):}
  \begin{itemize}
    \item \textit{Mục đích:} Bảo vệ máy chủ khỏi nguy cơ bị quét dò mật khẩu từ Internet (brute-force attack) và hỗ trợ chẩn đoán nhanh nguyên nhân khi không kết nối SSH được vào máy ảo.
    \item \textit{Quy trình người dùng:}
    \begin{itemize}
      \item Người dùng gửi yêu cầu mở cổng hoặc hỏi nguyên nhân sự cố (ví dụ: \textit{''Mở port 22 cho mọi người truy cập''} hoặc \textit{''Tôi không SSH được vào máy ảo i-0abc123, sao vậy?''}).
      \item \textit{Kiểm tra an toàn tường lửa:} Hệ thống từ chối mở port 22 cho toàn bộ dải mạng công khai (0.0.0.0/0), giải thích nguy cơ bị tấn công dò quét và gợi ý chỉ mở riêng cho địa chỉ IP hiện tại của người dùng.
      \item \textit{Chẩn đoán kết nối:} Nếu phát hiện người dùng mất file khóa SSH (.pem), hệ thống giải thích nguyên nhân và hướng dẫn sử dụng các kênh truy cập thay thế (như AWS Systems Manager Session Manager hoặc EC2 Instance Connect).
      \item Người dùng phản hồi đồng ý để áp dụng cấu hình an toàn được đề xuất.
    \end{itemize}
  \end{itemize}
\end{enumerate}

\subsubsection{Nhóm kịch bản và quy trình quản trị Amazon S3}

Dịch vụ lưu trữ đối tượng (Amazon Simple Storage Service - S3) dùng để lưu trữ tập tin, bản sao lưu và dữ liệu tĩnh của hệ thống. SatiBot hỗ trợ 06 quy trình quản lý S3 như sau:

\begin{enumerate}
  \item \textbf{Khởi tạo thùng chứa (S3 Bucket) an toàn và chuẩn hóa:}
  \begin{itemize}
    \item \textit{Mục đích:} Tạo kho lưu trữ đối tượng được thiết lập ở chế độ riêng tư theo mặc định và tuân thủ quy tắc đặt tên của AWS.
    \item \textit{Quy trình người dùng:} Nhắn yêu cầu tạo kèm tên mong muốn (ví dụ: \textit{''Tạo bucket tailieu-backup''}). Nếu tên bị trùng hoặc sai quy chuẩn đặt tên, xem hướng dẫn từ hệ thống và nhắn lại tên mới. Nhận thông báo bucket đã được tạo thành công ở chế độ riêng tư.
  \end{itemize}

  \item \textbf{Quét và khóa rủi ro lộ dữ liệu ra Internet (Security Scan \& Remediation):}
  \begin{itemize}
    \item \textit{Mục đích:} Rà soát và ngăn chặn kịp thời tình trạng lộ dữ liệu do vô tình mở quyền truy cập công khai (Public Access) trên bucket.
    \item \textit{Quy trình người dùng:} Nhắn yêu cầu kiểm tra (ví dụ: \textit{''Kiểm tra xem có bucket nào đang bị lộ ra ngoài Internet không?''}). Hệ thống quét chính sách truy cập của toàn bộ bucket trong tài khoản và trả về danh sách các bucket đang có nguy cơ. Người dùng xác nhận (ví dụ: \textit{''Khóa lại giùm tôi''}) để hệ thống bật tính năng Block Public Access ngay lập tức.
  \end{itemize}

  \item \textbf{Xóa bucket S3 an toàn có kiểm soát (chống xóa nhầm dữ liệu):}
  \begin{itemize}
    \item \textit{Mục đích:} Dọn dẹp bucket cũ mà không lo ngại nguy cơ xóa nhầm dữ liệu quan trọng đang lưu bên trong.
    \item \textit{Quy trình người dùng:} Người dùng nhắn yêu cầu xóa (ví dụ: \textit{''Xóa bucket du-an-2025''}). Hệ thống kiểm tra và thông báo số lượng tệp tin còn tồn bên trong. Người dùng phải nhập xác nhận rõ ràng; khi đó hệ thống mới tiến hành dọn sạch các đối tượng và xóa bucket.
  \end{itemize}

  \item \textbf{Tối ưu chi phí lưu trữ lâu dài (FinOps Lifecycle Rules):}
  \begin{itemize}
    \item \textit{Mục đích:} Giảm thiểu chi phí lưu trữ bằng cách tự động chuyển các tệp cũ, ít sử dụng sang các tầng lưu trữ giá rẻ hơn (S3 Standard-IA, S3 Glacier).
    \item \textit{Quy trình người dùng:} Người dùng nhắn yêu cầu chuyển tầng lưu trữ kèm mốc thời gian (ví dụ: \textit{''Tự chuyển file trong bucket sang kho lưu trữ giá rẻ sau 30 ngày''}). Hệ thống thiết lập quy tắc vòng đời (Lifecycle Configuration Rule) trên S3 và gửi thông báo xác nhận chính sách đã được áp dụng.
  \end{itemize}

  \item \textbf{Chia sẻ tệp tin an toàn có thời hạn (Presigned URL):}
  \begin{itemize}
    \item \textit{Mục đích:} Cung cấp liên kết tải tệp tin an toàn cho người ngoài trong khoảng thời gian nhất định mà không cần mở public dữ liệu hay cấp quyền tài khoản AWS.
    \item \textit{Quy trình người dùng:} Người dùng nhắn tên tệp và thời hạn chia sẻ (ví dụ: \textit{''Tạo link tải file hop-dong-2026.pdf trong 15 phút''}). Hệ thống tạo đường dẫn có chữ ký xác thực (Presigned URL) với thời gian hết hạn (TTL) 15 phút và gửi liên kết trực tiếp trên khung chat.
  \end{itemize}

  \item \textbf{Tra cứu dung lượng và tình trạng kho lưu trữ (Storage Analytics):}
  \begin{itemize}
    \item \textit{Mục đích:} Nắm bắt nhanh số lượng tệp tin, dung lượng sử dụng thực tế và dự toán chi phí lưu trữ của từng bucket.
    \item \textit{Quy trình người dùng:} Người dùng đặt câu hỏi tra cứu (ví dụ: \textit{''Bucket tailieu-duan có bao nhiêu file và chiếm bao nhiêu GB?''}). Hệ thống truy vấn CloudWatch Metrics và S3 API để thống kê chi tiết: số lượng objects, dung lượng tổng (GB), tầng lưu trữ và chi phí ước tính, hiển thị dưới dạng bảng rõ ràng.
  \end{itemize}
\end{enumerate}

\subsubsection{Nhóm kịch bản và quy trình quản trị AWS IAM}

Dịch vụ quản lý định danh và quyền truy cập (AWS Identity and Access Management - IAM) là trung tâm kiểm soát an ninh của hạ tầng AWS. SatiBot cung cấp 05 quy trình quản trị IAM an toàn, tuân thủ nguyên tắc đặc quyền tối thiểu (Least Privilege):

\begin{enumerate}
  \item \textbf{Tra cứu vai trò (Role) và chính sách quyền hạn (Policy):}
  \begin{itemize}
    \item \textit{Mục đích:} Xem nhanh danh sách hoặc kiểm tra chi tiết các quyền hạn đang được gán cho nhân sự hoặc dịch vụ để phục vụ công tác kiểm toán an ninh.
    \item \textit{Quy trình người dùng:} Người dùng nhắn câu hỏi tra cứu (ví dụ: \textit{''Kiểm tra xem Role lambda-s3-role đang có những quyền gì?''} hoặc \textit{''Liệt kê các Role dành cho dịch vụ EC2''}). Hệ thống phân tích tài liệu chính sách JSON và hiển thị bảng tóm tắt quyền hạn (tên quyền, dịch vụ được phép truy cập, phạm vi đọc/ghi) ngay trên khung chat.
  \end{itemize}

  \item \textbf{Khởi tạo vai trò (Role) hoặc chính sách (Policy) mới theo nguyên tắc tối thiểu:}
  \begin{itemize}
    \item \textit{Mục đích:} Tạo mới đối tượng phân quyền cho máy chủ, ứng dụng hoặc nhân sự mà không phải tự tay viết các đoạn mã JSON chính sách phức tạp và dễ nhầm lẫn.
    \item \textit{Quy trình người dùng:} Người dùng nhắn yêu cầu tạo kèm mục đích (ví dụ: \textit{''Tạo một Role cho máy ảo EC2 chỉ được phép đọc dữ liệu từ S3''} hoặc \textit{''Tạo Policy cho phép đọc log CloudWatch''}). Hệ thống đề xuất cấu trúc quyền hạn phù hợp (chỉ cấp quyền đọc dữ liệu, không cấp thừa quyền ghi/xóa). Sau khi người dùng xác nhận, hệ thống gọi API tạo đối tượng và trả về ARN của Role/Policy.
  \end{itemize}

  \item \textbf{Cấp hoặc điều chỉnh quyền truy cập:}
  \begin{itemize}
    \item \textit{Mục đích:} Bổ sung hoặc nới rộng quyền hạn cho nhân sự/dịch vụ khi phát sinh công việc mới trong dự án.
    \item \textit{Quy trình người dùng:} Người dùng nhắn yêu cầu gán thêm quyền (ví dụ: \textit{''Cấp thêm quyền đọc DynamoDB cho Role app-backend''}). Hệ thống đối soát đối tượng nhận quyền và chính sách tương ứng, yêu cầu người dùng xác nhận và trả về thông báo quyền truy cập mới đã được áp dụng thành công.
  \end{itemize}

  \item \textbf{Thu hồi quyền truy cập (Least Privilege Enforcement):}
  \begin{itemize}
    \item \textit{Mục đích:} Gỡ bỏ các quyền hạn không còn sử dụng hoặc thu hồi các quyền nguy hiểm (như quyền AdministratorAccess cấp thừa cho tài khoản thực tập sinh) nhằm hạn chế rủi ro an toàn thông tin.
    \item \textit{Quy trình người dùng:} Người dùng nhắn yêu cầu gỡ bỏ quyền (ví dụ: \textit{''Gỡ quyền AdministratorAccess khỏi tài khoản dev-intern''}). Hệ thống hiển thị thông báo phạm vi quyền chuẩn bị gỡ bỏ; sau khi người dùng xác nhận, hệ thống gỡ bỏ chính sách và thông báo kết quả.
  \end{itemize}

  \item \textbf{Xóa đối tượng IAM an toàn:}
  \begin{itemize}
    \item \textit{Mục đích:} Dọn dẹp dứt điểm các Role, Policy hoặc User cũ không còn sử dụng để tránh dư thừa tài nguyên và loại bỏ lỗ hổng bảo mật tiềm tàng.
    \item \textit{Quy trình người dùng:} Người dùng gửi yêu cầu xóa (ví dụ: \textit{''Xóa Role test-migration cũ giùm tôi''}). Hệ thống kiểm tra các quan hệ phụ thuộc (nếu Role đang gắn với máy ảo EC2 đang chạy, hệ thống sẽ cảnh báo nguy cơ làm dừng ứng dụng). Người dùng xác nhận lần cuối để hệ thống tự động tháo gỡ các liên kết an toàn trước khi xóa bỏ đối tượng IAM.
  \end{itemize}
\end{enumerate}

\subsubsection{Nhóm kịch bản và quy trình quản trị AWS Lambda}

Dịch vụ điện toán không máy chủ (AWS Lambda) cho phép chạy mã nguồn theo sự kiện. SatiBot hỗ trợ 05 quy trình quản trị và kiểm thử hàm serverless trực tiếp:

\begin{enumerate}
  \item \textbf{Tra cứu danh mục và kiểm tra cấu hình hàm (Inventory \& Config Audit):}
  \begin{itemize}
    \item \textit{Mục đích:} Nắm bắt nhanh toàn bộ các hàm Lambda đang có trong tài khoản cùng các thông số vận hành cốt lõi (ngôn ngữ runtime, bộ nhớ RAM, thời gian timeout) phục vụ công tác kiểm toán tài nguyên.
    \item \textit{Quy trình người dùng:} Người dùng gửi câu hỏi (ví dụ: \textit{''Liệt kê các hàm Lambda trong tài khoản''} hoặc \textit{''Hàm image-processor đang chạy ngôn ngữ gì và cấp bao nhiêu RAM?''}). Hệ thống hiển thị bảng tổng hợp gồm Tên hàm, Môi trường thực thi (Runtime), Dung lượng bộ nhớ (Memory) và Thời gian chờ tối đa (Timeout).
  \end{itemize}

  \item \textbf{Chạy thử nghiệm hàm trực tiếp qua Chat (Direct Invocation):}
  \begin{itemize}
    \item \textit{Mục đích:} Kích hoạt và kiểm tra nhanh kết quả trả về của hàm với dữ liệu thử nghiệm ngay trên khung chat mà không cần đợi sự kiện bên ngoài hay cấu hình test event trên Console.
    \item \textit{Quy trình người dùng:} Người dùng nhắn lệnh chạy kèm dữ liệu đầu vào (ví dụ: \textit{''Chạy thử hàm calculate-price với dữ liệu giá 100 và thuế 10''}). Hệ thống đóng gói dữ liệu JSON, gọi API `invoke`, đo thời gian thực thi và trả về kết quả (Payload phản hồi, mã HTTP status, thời gian chạy tính bằng mili-giây) trực tiếp trên khung chat.
  \end{itemize}

  \item \textbf{Điều chỉnh cấu hình bộ nhớ và thời gian chờ của hàm Lambda:}
  \begin{itemize}
    \item \textit{Mục đích:} Nâng cấp cấu hình tức thời cho các hàm đang chạy chậm hoặc có nguy cơ sập do thiếu bộ nhớ hay hết thời gian chờ mà không phải can thiệp mã nguồn hay mở Console.
    \item \textit{Quy trình người dùng:} Người dùng ra lệnh điều chỉnh thông số (ví dụ: \textit{''Tăng thời gian timeout của hàm backup-data lên 30 giây''} hoặc \textit{''Nâng RAM của hàm image-processor lên 512 MB''}). Hệ thống hiển thị bảng đối soát thông số cũ và thông số mới. Sau khi người dùng xác nhận, hệ thống gọi API cập nhật cấu hình và thông báo thành công.
  \end{itemize}

  \item \textbf{Cấp phát nhanh liên kết HTTP công khai (Function URL Generation):}
  \begin{itemize}
    \item \textit{Mục đích:} Tạo nhanh một đường dẫn Endpoint HTTPS công khai (Function URL) để người dùng hoặc webhook bên ngoài có thể gọi trực tiếp hàm Lambda mà không mất công cấu hình API Gateway phức tạp.
    \item \textit{Quy trình người dùng:} Người dùng nhắn: \textit{''Tạo link Function URL công khai cho hàm image-api''}. Hệ thống tự động cấu hình Function URL, thiết lập quyền truy cập và gửi lại đường link HTTPS dạng \url{https://<url-id>.lambda-url.<region>.on.aws/} trực tiếp trên khung chat để người dùng kiểm chứng kết quả thực thi của hàm.
  \end{itemize}

  \item \textbf{Tạm dừng hoạt động hàm khẩn cấp (Emergency Kill-Switch):}
  \begin{itemize}
    \item \textit{Mục đích:} Vô hiệu hóa ngay lập tức một hàm đang bị lặp lỗi logic gây tiêu hao ngân sách hoặc cần ngắt kết nối gấp để bảo trì hệ thống mà không cần xóa bỏ mã nguồn của hàm.
    \item \textit{Quy trình người dùng:} Người dùng nhắn lệnh ngắt khẩn cấp (ví dụ: \textit{''Khóa tạm thời hàm order-service để bảo trì hệ thống''}). Hệ thống cảnh báo về việc hàm sẽ từ chối mọi yêu cầu mới (trả về lỗi HTTP 429 Too Many Requests) và đề nghị người dùng xác nhận. Sau khi nhận xác nhận, hệ thống đặt mức thực thi đồng thời về 0 (\texttt{Reserved Concurrency = 0}), đưa hàm vào trạng thái tạm ngừng hoạt động an toàn. Người dùng có thể yêu cầu mở lại mức concurrency bất cứ lúc nào.
  \end{itemize}
\end{enumerate}
\endgroup

\clearpage

\section{GIẢI PHÁP THỰC HIỆN DỰ KIẾN}

\subsection{Kiến trúc tính năng của hệ thống}

Hình~\ref{fig:kientruc_tinhnang} mô tả bốn phân hệ chính của SatiBot và quan hệ giữa chúng: xác thực người dùng, kết nối tài khoản AWS, quản trị tài nguyên qua Chatbot và tra cứu lịch sử thao tác. Luồng nghiệp vụ tổng quát là: người dùng đăng nhập và liên kết tài khoản AWS, gửi yêu cầu bằng tiếng Việt, xem phương án hệ thống đề xuất, xác nhận thao tác, rồi nhận kết quả thực thi.

\begin{figure}[H]
\centering
\IfFileExists{Image/So_do_kien_truc_tinh_nang.jpg}{
  \includegraphics[width=0.74\textwidth]{Image/So_do_kien_truc_tinh_nang.jpg}
}{
  \fbox{\parbox[c][6.5cm][c]{0.88\textwidth}{\centering Sơ đồ kiến trúc tính năng hệ thống SatiBot}}
}
\caption{Sơ đồ kiến trúc tính năng và luồng xử lý nghiệp vụ của SatiBot}
\label{fig:kientruc_tinhnang}
\end{figure}

Các phân hệ phối hợp như sau:
\begin{enumerate}[itemsep=3pt,topsep=2pt]
  \item \textbf{Đăng ký và xác thực người dùng:} Amazon Cognito xử lý đăng ký, đăng nhập và cấp JWT để xác thực các yêu cầu nghiệp vụ. Sau khi đăng ký, Lambda tạo hồ sơ người dùng trong bảng \texttt{Users} của DynamoDB.

  \item \textbf{Kết nối tài khoản AWS:} SatiBot tạo \texttt{ExternalId}; người dùng dùng mã này và hướng dẫn của hệ thống để tạo IAM Role trong tài khoản AWS của mình, sau đó gửi \texttt{RoleArn} về SatiBot. Connector Lambda gọi STS \texttt{AssumeRole} để kiểm tra khả năng ủy quyền. Nếu thành công, hệ thống lưu thông tin kết nối và trạng thái trong bảng \texttt{AwsConnections}; SatiBot không cần lưu Access Key/Secret Key dài hạn của người dùng. Luồng chi tiết được minh họa ở Hình~\ref{fig:sequence_connect}.


  \begin{figure}[H]
  \centering
  \IfFileExists{Image/sequence3.jpg}{
    \includegraphics[width=0.78\textwidth]{Image/sequence3.jpg}
  }{
    \fbox{\parbox[c][5cm][c]{0.88\textwidth}{\centering Sơ đồ tuần tự quy trình kết nối tài khoản AWS}}
  }
  \caption{Sơ đồ tuần tự quy trình kết nối và xác thực tài khoản AWS mục tiêu}
  \label{fig:sequence_connect}
  \end{figure}

  \item \textbf{Quản trị tài nguyên qua Chatbot:} Chatbot Lambda nhận câu lệnh, lấy ngữ cảnh hội thoại từ \texttt{SessionState} và gửi yêu cầu cùng Function Calling Schema đến mô hình LLM. Kết quả được kiểm tra về tham số và quy tắc an toàn trước khi lưu thành hành động chờ xác nhận trong \texttt{PendingAction}. Khi người dùng xác nhận, hệ thống kiểm tra hành động còn hiệu lực và thuộc đúng người dùng; Dispatcher Lambda dùng STS và Boto3 để gọi API của EC2, S3, IAM hoặc Lambda. Kết quả được lưu vào \texttt{ActionHistory} và gửi về giao diện qua AppSync. Luồng chi tiết được minh họa ở Hình~\ref{fig:sequence_chatops}.

  \begin{figure}[H]
  \centering
  \IfFileExists{Image/sequence1.jpg}{
    \includegraphics[width=0.58\textwidth]{Image/sequence1.jpg}
  }{
    \fbox{\parbox[c][6.5cm][c]{0.88\textwidth}{\centering Sơ đồ tuần tự quy trình quản trị tài nguyên qua Chatbot}}
  }
  \caption{Sơ đồ tuần tự quy trình điều phối và thực thi tác vụ ChatOps hai giai đoạn}
  \label{fig:sequence_chatops}
  \end{figure}

  \item \textbf{Tra cứu lịch sử thao tác:} History Lambda truy vấn bảng \texttt{ActionHistory} theo người dùng và các điều kiện tra cứu, sau đó trả về thời gian, hành động, tài nguyên đích và trạng thái thực thi để hiển thị trên giao diện.
\end{enumerate}

\subsection{Kiến trúc triển khai của hệ thống}

SatiBot sử dụng các dịch vụ AWS được quản lý để triển khai giao diện, xác thực, xử lý nghiệp vụ và lưu trữ dữ liệu. Hình~\ref{fig:so_do_kien_truc} thể hiện cách các thành phần kết nối; Bảng~\ref{tab:kientruc} tóm tắt trách nhiệm của từng dịch vụ.

\begin{table}[H]
\centering
\caption{Các thành phần và dịch vụ trong kiến trúc triển khai hệ thống SatiBot}
\label{tab:kientruc}
\bangfont
\begin{tabularx}{\linewidth}{@{}>{\raggedright\arraybackslash}p{4.5cm}Y@{}}
\toprule
\textbf{Thành phần / Dịch vụ} & \textbf{Vai trò trong kiến trúc}\\
\midrule
React, Amazon S3 và CloudFront & Cung cấp giao diện ChatOps; S3 lưu các tệp tĩnh và CloudFront phân phối giao diện qua HTTPS.\\
Amazon Cognito & Quản lý đăng ký, đăng nhập và cấp JWT để xác thực người dùng.\\
Amazon API Gateway & Tiếp nhận API từ frontend, xác thực JWT và định tuyến yêu cầu chat, xác nhận thao tác hoặc tra cứu lịch sử.\\
AWS Step Functions & Điều phối luồng thực thi sau khi người dùng xác nhận, gồm Dispatcher và Formatter Lambda.\\
AWS Lambda & Thực hiện các nghiệp vụ: Chatbot, Dispatcher, Formatter, Connector, History và tạo External ID.\\
Mô hình LLM & Phân tích câu lệnh tiếng Việt và trả về ý định, tham số theo Function Calling Schema; không trực tiếp gọi API AWS.\\
AWS STS và Boto3 & STS cấp thông tin xác thực tạm thời cho tài khoản đích; Boto3 dùng thông tin này để gọi API tài nguyên.\\
Amazon DynamoDB & Lưu hồ sơ người dùng, kết nối AWS, lịch sử hội thoại, hành động chờ xác nhận và lịch sử thao tác.\\
AWS AppSync & Phát kết quả qua GraphQL Subscriptions để giao diện cập nhật mà không cần polling liên tục.\\
Amazon CloudWatch & Lưu log thực thi và hỗ trợ theo dõi, chẩn đoán lỗi.\\
\bottomrule
\end{tabularx}
\end{table}

\subsubsection{Sơ đồ kiến trúc triển khai tổng thể và giao diện tương tác}

\begin{figure}[H]
\centering
\IfFileExists{Image/so_do_kien_truc.jpg}{
  \includegraphics[width=0.82\textwidth]{Image/so_do_kien_truc.jpg}
}{
  \fbox{\parbox[c][6.5cm][c]{0.88\textwidth}{\centering Sơ đồ kiến trúc triển khai tổng thể hệ thống SatiBot}}
}
\caption{Sơ đồ kiến trúc triển khai tổng thể hệ thống SatiBot}
\label{fig:so_do_kien_truc}
\end{figure}

\begin{figure}[H]
\centering
\IfFileExists{Image/so_do_erd.png}{
  \includegraphics[width=0.62\textwidth]{Image/so_do_erd.png}
}{
  \fbox{\parbox[c][6.5cm][c]{0.88\textwidth}{\centering Sơ đồ thực thể liên kết ERD}}
}
\caption{Sơ đồ thực thể liên kết (ERD) cơ sở dữ liệu hệ thống}
\label{fig:so_do_erd}
\end{figure}

Luồng hoạt động và giao tiếp giữa các thành phần gồm bốn bước:
\begin{enumerate}[itemsep=3pt,topsep=2pt]
  \item \textbf{Tiếp nhận và phân tích yêu cầu:} Người dùng gửi câu lệnh đến API Gateway qua HTTPS. Sau khi kiểm tra JWT, Gateway chuyển yêu cầu đến Chatbot Lambda. Hàm này lấy ngữ cảnh hội thoại từ DynamoDB, gọi mô hình LLM để xác định hành động và tham số, rồi kiểm tra đầu ra theo schema và quy tắc an toàn.

  \item \textbf{Tạo yêu cầu chờ xác nhận:} Nếu thông tin hợp lệ, hệ thống lưu phương án vào \texttt{PendingAction} và trả phần tóm tắt cho frontend. Nếu thiếu tham số hoặc có rủi ro, hệ thống yêu cầu người dùng bổ sung thông tin hoặc điều chỉnh phương án; chưa thực thi thao tác làm thay đổi tài nguyên.

  \item \textbf{Xác nhận và thực thi:} Khi người dùng xác nhận, frontend gửi mã hành động đến API Gateway. Hệ thống kiểm tra quyền sở hữu, trạng thái và thời hạn của hành động rồi khởi chạy Step Functions. Dispatcher Lambda lấy thông tin kết nối từ \texttt{AwsConnections}, gọi STS \texttt{AssumeRole} với \texttt{RoleArn} và \texttt{ExternalId}, sau đó dùng Boto3 thực hiện thao tác trên tài khoản AWS đích.

  \item \textbf{Ghi nhận và trả kết quả:} Kết quả thành công hoặc thất bại được Formatter Lambda chuẩn hóa và ghi vào \texttt{ActionHistory}. Formatter gửi thông báo qua GraphQL Mutation đến AppSync; frontend nhận kết quả qua GraphQL Subscription. Log kỹ thuật được ghi vào CloudWatch để phục vụ kiểm tra lỗi.
\end{enumerate}

Bảng~\ref{tab:inter_service_interfaces} tóm tắt các giao diện chính, giao thức, cơ chế xác thực và dữ liệu trao đổi giữa các thành phần.

\begin{table}[H]
\centering
\caption{Đặc tả giao diện và giao thức tương tác giữa các thành phần}
\label{tab:inter_service_interfaces}
\begingroup
\renewcommand{\arraystretch}{1.12}
\fontsize{8.5}{11.2}\selectfont
\begin{tabularx}{\linewidth}{@{}>{\raggedright\arraybackslash}p{4.1cm}>{\raggedright\arraybackslash}p{2.35cm}>{\raggedright\arraybackslash}p{3.05cm}Y@{}}
\toprule
\textbf{Nguồn $\rightarrow$ Đích} & \textbf{Giao thức} & \textbf{Xác thực} & \textbf{Dữ liệu trao đổi}\\
\midrule
Web Client $\rightarrow$ CloudFront $\rightarrow$ S3 & HTTPS & TLS & Tệp giao diện tĩnh\\
Web Client $\rightarrow$ API Gateway & HTTPS REST & Cognito JWT & Câu lệnh, session ID, mã hành động\\
API Gateway $\rightarrow$ Chatbot Lambda & Lambda Proxy & IAM & Nội dung chat và thông tin phiên\\
API Gateway $\rightarrow$ Step Functions & StartExecution & IAM Role & User ID và mã hành động đã xác nhận\\
Chatbot Lambda $\leftrightarrow$ LLM & HTTPS API & Theo nhà cung cấp LLM & Prompt, schema công cụ, tham số và kết quả phân tích\\
Step Functions $\rightarrow$ Lambda Tasks & AWS SDK Integration & Step Functions Execution Role & JSON đầu vào/đầu ra từng bước\\
Lambda $\leftrightarrow$ DynamoDB & Boto3 SDK & Lambda Execution Role & GetItem, PutItem, Query, UpdateItem\\
Dispatcher Lambda $\rightarrow$ AWS STS & Boto3 / HTTPS & Lambda Execution Role & AssumeRole với RoleArn và ExternalId\\
Dispatcher $\rightarrow$ Tài khoản AWS đích & Boto3 / HTTPS & Credential tạm thời, SigV4 & Yêu cầu API và kết quả thao tác AWS\\
Formatter Lambda $\rightarrow$ AppSync & GraphQL Mutation / HTTPS & IAM SigV4 & Nội dung kết quả cần phát tới frontend\\
Web Client $\leftrightarrow$ AppSync & GraphQL Subscription / WSS & Cognito User Pool & Sự kiện cập nhật trên khung chat\\
\bottomrule
\end{tabularx}
\endgroup
\end{table}

\clearpage

\section{KẾ HOẠCH TRIỂN KHAI DỰ ÁN}

\subsection{Kế hoạch triển khai theo từng giai đoạn (05 tuần)}

Dự án được lên kế hoạch thực hiện trong 05 tuần theo mô hình Agile/Scrum, chia thành 05 Sprint tương ứng với từng tuần phát triển. Mỗi tuần tập trung hoàn thành một nhóm mục tiêu cụ thể, kiểm thử độc lập từng phần trước khi ghép nối hệ thống:

\begin{itemize}
  \item \textbf{Tuần 1 (Sprint 1) -- Chuẩn bị môi trường và kiểm thử PoC độc lập:} Thống nhất quy định làm việc, quy chuẩn viết code chung. Tạo tài khoản AWS chung và IAM User cho từng thành viên. Mỗi thành viên tìm hiểu tài liệu kỹ thuật, viết và kiểm thử độc lập các hàm Python/Boto3 cơ bản cho dịch vụ AWS mình phụ trách (S3, EC2, Lambda, IAM).
  \item \textbf{Tuần 2 (Sprint 2) -- Xây dựng hạ tầng cơ bản và ghép nối Backend:} Dựng khung giao diện web React tĩnh trên S3/CloudFront; cấu hình đăng ký, đăng nhập tài khoản qua Amazon Cognito. Thiết lập API Gateway, viết Lambda Backend tiếp nhận yêu cầu, tạo 05 bảng trên DynamoDB và tích hợp các module Boto3 của 4 thành viên vào Backend.
  \item \textbf{Tuần 3 (Sprint 3) -- Tích hợp AI Chatbot, Step Functions và AppSync:} Tích hợp mô hình ngôn ngữ lớn (LLM), viết prompt và khai báo cấu trúc Function Calling (Tool Schemas). Dựng máy trạng thái Step Functions để điều phối luồng: nhận câu chat $\rightarrow$ Backend gọi LLM trích xuất tham số $\rightarrow$ hiển thị xác nhận $\rightarrow$ xác thực qua STS AssumeRole $\rightarrow$ Boto3 thực thi trên AWS $\rightarrow$ trả kết quả lên giao diện qua kênh WebSocket của AWS AppSync.
  \item \textbf{Tuần 4 (Sprint 4) -- Kiểm thử hệ thống và tối ưu hóa:} Kiểm thử các kịch bản thực tế, bao gồm cả trường hợp chạy thành công và các tình huống thiếu quyền (AccessDenied); cấu hình CloudWatch Logs để theo dõi lỗi tập trung; tối ưu thời gian phản hồi của Lambda và hoàn thiện giao diện chat.
  \item \textbf{Tuần 5 (Sprint 5) -- Nghiệm thu, dọn dẹp tài nguyên và hoàn thiện báo cáo:} Chạy thử toàn bộ các ca sử dụng để nghiệm thu sản phẩm; hoàn thiện báo cáo kỹ thuật, slide thuyết trình; rà soát và xóa các tài nguyên AWS thử nghiệm để tránh phát sinh chi phí.
\end{itemize}

\subsection{Kế hoạch hành động chi tiết theo từng Sprint}

Bảng~\ref{tab:action_plan} liệt kê chi tiết từng đầu việc trong kế hoạch hành động (Action Plan), mối liên hệ với các thành phần trong kiến trúc hệ thống và phân công trách nhiệm cho 04 thành viên tương ứng với 05 tuần (05 Sprint) thực hiện dự án.

\begingroup
\setlength{\tabcolsep}{3.5pt}
\renewcommand{\arraystretch}{1.0}
\fontsize{8.5}{10.8}\selectfont
\begin{longtable}{@{}p{0.6cm}>{\raggedright\arraybackslash}p{3.8cm}>{\raggedright\arraybackslash}p{3.2cm}>{\raggedright\arraybackslash}p{2.6cm}>{\raggedright\arraybackslash}p{4.8cm}@{}}
\caption{Kế hoạch hành động và phân công công việc chi tiết theo từng Sprint}\label{tab:action_plan}\\
\toprule
\textbf{STT} & \textbf{Hạng mục công việc} & \textbf{Thành phần kiến trúc liên quan} & \textbf{Người phụ trách} & \textbf{Kết quả bàn giao}\\
\midrule
\endfirsthead
\caption[]{(Tiếp theo Bảng~\ref{tab:action_plan})}\\
\toprule
\textbf{STT} & \textbf{Hạng mục công việc} & \textbf{Thành phần kiến trúc liên quan} & \textbf{Người phụ trách} & \textbf{Kết quả bàn giao}\\
\midrule
\endhead
\midrule
\multicolumn{5}{r}{\textit{(Còn tiếp sang trang sau...)}}\\
\endfoot
\bottomrule
\endlastfoot
\multicolumn{5}{@{}l}{\textbf{Giai đoạn 1: Tuần 1 -- Thiết lập quy trình, hạ tầng ban đầu và PoC Boto3 độc lập}}\\
1 & Viết và kiểm thử độc lập module Boto3 cho dịch vụ S3 & User's Cloud (Amazon S3) & Phạm Đức Hùng & Module \texttt{s3\_service} (tạo bucket, presign, block public).\\
2 & Thống nhất kịch bản và kiểm thử module Boto3 cho dịch vụ EC2 & User's Cloud (Amazon EC2) & Nguyễn Đức Huy & Module \texttt{ec2\_service} (khởi tạo, bật/tắt, describe).\\
3 & Thống nhất kịch bản và kiểm thử module Boto3 cho Lambda & User's Cloud (AWS Lambda) & Trịnh Thiên Quang & Module \texttt{lambda\_service} (list, invoke, hot config).\\
4 & Thống nhất quy tắc làm việc, tạo tài khoản AWS và module IAM & Toàn bộ kiến trúc, AWS IAM & Cao Anh Minh & Quy chuẩn mã nguồn, AWS Account chung, module \texttt{iam\_service}.\\
\midrule
\multicolumn{5}{@{}l}{\textbf{Giai đoạn 2: Tuần 2 -- Dựng hạ tầng nền tảng Serverless, Giao diện và Tích hợp Backend}}\\
5 & Phát triển khung giao diện React cơ bản trên S3/CloudFront & Frontend (React, S3, CloudFront) & Phạm Đức Hùng & Giao diện Web Chat responsive, CDN CloudFront với HTTPS.\\
6 & Cấu hình luồng đăng nhập và quản lý danh tính qua Cognito & Amazon Cognito, API Gateway & Nguyễn Đức Huy & Luồng Login/Register qua Cognito, JWT Authorizer trên Gateway.\\
7 & Dựng Lambda Backend trung tâm và tích hợp các module Boto3 & AWS Lambda Backend & Trịnh Thiên Quang & Lambda Backend tiếp nhận request và tích hợp các module Boto3.\\
8 & Cấu hình IAM Role cho Backend và thiết kế 05 bảng DynamoDB & AWS IAM, DynamoDB & Cao Anh Minh & IAM Execution Role, 05 bảng DynamoDB cơ sở.\\
\midrule
\multicolumn{5}{@{}l}{\textbf{Giai đoạn 3: Tuần 3 -- Tích hợp AI Chatbot, Step Functions và AppSync}}\\
9 & Định nghĩa Function Calling và Prompt cho dịch vụ S3 & Model LLM, Chatbot Lambda & Phạm Đức Hùng & Bộ Tool Schema cho S3, hỗ trợ kênh WebSocket AppSync.\\
10 & Định nghĩa Function Calling cho EC2 và ghép luồng LLM & Model LLM, Chatbot Lambda & Nguyễn Đức Huy & Bộ Tool Schema cho EC2, hoàn thiện ghép luồng LLM.\\
11 & Tích hợp API LLM, xử lý phản hồi Function Calling tại Backend & Model LLM, Lambda Backend & Trịnh Thiên Quang & Xử lý phản hồi Function Calling và tích hợp LLM vào Backend.\\
12 & Hoàn thiện Prompt IAM, xây dựng Step Functions State Machine & AWS Step Functions, STS & Cao Anh Minh & State Machine điều phối luồng thực thi, xác thực STS.\\
\midrule
\multicolumn{5}{@{}l}{\textbf{Giai đoạn 4: Tuần 4 -- Kiểm thử toàn diện, Phân quyền và Tối ưu hóa}}\\
13 & Kiểm thử hệ thống, tối ưu hóa module S3 & Phân hệ S3, Web Client & Phạm Đức Hùng & Báo cáo kiểm thử S3, tối ưu giao diện truyền tải file.\\
14 & Kiểm thử hệ thống, tối ưu module EC2, dọn dẹp tài nguyên & Phân hệ EC2, CloudWatch & Nguyễn Đức Huy & Báo cáo kiểm thử EC2, kịch bản FinOps cảnh báo máy ảo.\\
15 & Kiểm thử phân quyền, tối ưu hệ thống Backend & Backend Lambda, Step Functions & Trịnh Thiên Quang & Báo cáo kiểm thử phân quyền, tối ưu Cold-start của Lambda.\\
16 & Kiểm thử trường hợp thiếu quyền (AccessDenied), CloudWatch Log & AWS IAM, CloudWatch Logs & Cao Anh Minh & Xử lý lỗi phân quyền, tập trung hóa log trên CloudWatch.\\
\midrule
\multicolumn{5}{@{}l}{\textbf{Giai đoạn 5: Tuần 5 -- Nghiệm thu, Tổng kết dự án và Đóng gói báo cáo}}\\
17 & Hoàn thiện báo cáo phần S3, hỗ trợ làm slide thuyết trình & Báo cáo đề tài & Phạm Đức Hùng & Nội dung báo cáo S3, hoàn thiện slide demo.\\
18 & Tổng hợp kết quả nghiệm thu EC2, chuẩn bị slide demo & Báo cáo đề tài & Nguyễn Đức Huy & Nội dung báo cáo EC2, slide thuyết trình kỹ thuật.\\
19 & Nghiệm thu luồng Lambda, hoàn thiện tài liệu hướng dẫn & Báo cáo đề tài & Trịnh Thiên Quang & Nội dung báo cáo Lambda, tài liệu vận hành.\\
20 & Tổng kết dự án, hoàn thiện báo cáo tổng thể, slide thuyết trình & Báo cáo đề tài & Cao Anh Minh & Bản báo cáo hoàn chỉnh (\texttt{report1.pdf}), slide tổng thể.\\
\end{longtable}
\endgroup

\end{document}